\documentclass[a4paper]{article}

\usepackage{amsmath}
\usepackage{enumitem}
\usepackage{hyperref}
\usepackage{xcolor}

\newcommand{\entropy}[1]{\mathcal{H}(#1)}
\newcommand{\fnin}[2]{\((#1, #2)\)}

\newcommand{\cmnt}[2]{{\color{red}{\textbf{#1:} #2}}}
\newcommand{\AL}[1]{\cmnt{AL}{#1}}

\begin{document}

Suppose we have a binary branch in a deterministic, self-contained (no external
inputs) piece of code, similar to an \texttt{if} statement present in most, if
not all, programming languages. In this document, we aim to compute how does
this affect the Uncertainty Coefficient~\cite{press2007numerical} in terms of
information flowing through either branches.

The formula for computing the uncertainty coefficient formula is given
in~\autoref{eq:unc_coef}. As the code is deterministic, we can further simplify
this formula by using Squeeziness~\cite{clark2012squeeziness}, as
per~\autoref{eq:squeez}. The final equation we use is given
in~\autoref{eq:uc_simple}.

\begin{equation}\label{eq:unc_coef}
    U(I|O) = \frac{\entropy{I} - \entropy{I|O}}{\entropy{I}} = 1 - \frac{\entropy{I|O}}{\entropy{I}}
\end{equation}

\begin{equation}\label{eq:squeez}
    \entropy{I|O} = \entropy{I} - \entropy{O}
\end{equation}

\begin{equation}\label{eq:uc_simple}
    U(I|O) = 1 - \frac{\entropy{I|O}}{\entropy{I}} = 1 - \frac{\entropy{I} - \entropy{O}}{\entropy{I}} = \frac{\entropy{O}}{\entropy{I}}
\end{equation}

Suppose at the entry of the conditional branch, there are \(\mathcal{U}^N\)
states, equivalent to \(N\) distinct states. Further suppose the condition in
the branch has a probability of \(p\) of being \texttt{true} (and thus a
probability of \(1-p\) of being \texttt{false}). Thus, out of the
\(\mathcal{U}^N\) total possible states entering the conditional,
\(p*\mathcal{U}^N\) will follow the \texttt{true} branch, with \((1-p) *
\mathcal{U}^N\) following the \texttt{false} branch.

\bibliographystyle{plain}
\bibliography{notes}

\end{document}
